\section{Superelliptic conductor basis and reference assembly algorithm}
\label{sec:superellipse-basis}

This section specifies a first implementable basis family for the homogeneous-
background conductor teacher. It is not the only admissible basis, but it is
the reference construction for the first implementation.

\subsection{Reference superellipse}

Let the conductor cross-section be
\[
 \Sigma
 =
 \left\{
  (\xi_1,\xi_2):
  \left|\frac{\xi_1}{a_c}\right|^m
  +
  \left|\frac{\xi_2}{b_c}\right|^m
  \le 1
 \right\},
 \qquad
 a_c=w/2,\quad b_c=h/2 .
\]
Use generalized polar coordinates
\[
 \xi_1
 =
 a_c\,\rho\,
 \operatorname{sgn}(\cos\theta)
 |\cos\theta|^{2/m},
\]
\[
 \xi_2
 =
 b_c\,\rho\,
 \operatorname{sgn}(\sin\theta)
 |\sin\theta|^{2/m},
\]
with
\[
 0\le\rho\le1,\qquad 0\le\theta<2\pi.
\]
The Jacobian is evaluated analytically or by automatic differentiation and is
included in all cross-section inner products. Since this coordinate chart has
integrable axis singularities for some exponents, production quadrature may
instead split the superellipse into four smooth quadrant patches. The
mathematical basis is independent of that quadrature choice.

\subsection{Scalar cross-section modes}

Define weighted radial functions
\[
 R_n^{(\beta)}(\rho)
 =
 (1-\rho)^{\beta/2}
 P_n^{(\alpha,\beta)}(2\rho-1),
\]
where $P_n^{(\alpha,\beta)}$ is a Jacobi polynomial. The boundary exponent
$\beta$ may be chosen to concentrate resolution near the conductor boundary
when skin depth is small.

Angular functions are
\[
 A_0(\theta)=1,
\]
\[
 A_{k,c}(\theta)=\cos(k\theta),\qquad
 A_{k,s}(\theta)=\sin(k\theta).
\]
A raw scalar family is
\[
 \varphi_{n,k,\tau}(\rho,\theta)
 =
 R_n^{(\beta_k)}(\rho)A_{k,\tau}(\theta).
\]
These functions are Gram--Schmidt or Cholesky orthonormalized under the
conductive cross-section metric
\[
 M_\Sigma(u,v)
 =
 \int_{\Sigma}u^\ast \sigma^{-1}v\,\dd A .
\]
The constant mode is retained exactly.

\subsection{Boundary-layer enrichment}

When
\[
 \chi_{\min}
 =
 \min(w,h)/\delta
\]
is large, a polynomial radial family may converge slowly. Add boundary-layer
modes
\[
 B_{r,k,\tau}
 =
 \exp[-\lambda_r(1-\rho)]
 A_{k,\tau}(\theta),
\]
with
\[
 \lambda_r
 \sim
 c_r\,\frac{\min(w,h)}{\delta}.
\]
The coefficients $c_r$ are chosen from a small geometric sequence. These
functions represent skin concentration without shrinking a surrounding volume
mesh.

\subsection{Longitudinal basis}

Parameterize the centerline by normalized arc coordinate
\[
 \zeta=s/L\in[0,1].
\]
The first implementation uses compact B-spline functions
\[
 \chi_\ell(\zeta)
\]
of degree $p_s\ge2$ with open knots. Knot spacing is refined according to:
\begin{enumerate}
\item centerline curvature;
\item terminal neighborhoods;
\item rapid variation of cross-section;
\item close electromagnetic proximity to another object.
\end{enumerate}

The longitudinal current basis is
\[
 \mathbf w_{\ell a}^{(t)}
 =
 e_t(s)\chi_\ell(s)\varphi_a(\xi;s).
\]

\subsection{Transverse compatible modes}

To allow finite-cross-section charge redistribution and transverse proximity
currents, define cross-section scalar potentials $\eta_a$ and stream functions
$\psi_a$. The transverse family is
\[
 \mathbf v_a^{(g)}
 =
 \nabla_\perp \eta_a,
\qquad
 \mathbf v_a^{(c)}
 =
 e_t\times\nabla_\perp\psi_a.
\]
The gradient family contributes cross-section divergence; the rotated-gradient
family is divergence-free in the local plane. Mapped to the swept conductor,
metric/curvature corrections are included through the exact geometry
Jacobian. The discrete divergence matrix $D$ is assembled from the mapped
basis, rather than assuming the straight-conductor formula.

\subsection{Charge basis}

Use longitudinal B-splines times scalar cross-section modes:
\[
 \nu_{\ell a}(s,\xi)
 =
 \widehat\chi_\ell(s)\widehat\varphi_a(\xi;s).
\]
The charge space is selected so the image of the discrete current divergence is
represented to machine precision or to a declared commuting-projection
tolerance. In the first implementation the safest policy is to build $D$ first,
compute its numerical range, and retain a charge basis spanning that range plus
terminal charge modes.

\subsection{Port terminal basis}

For each terminal cross-section, define a normalized flux mode
\[
 \tau_p(\xi)
\]
such that
\[
 \int_{\Gamma_p}\tau_p(\xi)\dd A=1.
\]
The terminal current constraint is imposed exactly:
\[
 \int_{\Gamma_p}\mathbf J\cdot n\,\dd A=I_p.
\]
Additional zero-net-flux terminal shape modes may be included to represent
nonuniform terminal current without changing the prescribed port current.

\subsection{Matrix elements}

For current basis functions $\mathbf w_\alpha,\mathbf w_\beta$,
\[
 R_{\alpha\beta}
 =
 \int_{\Omega_c}
 \mathbf w_\alpha^\ast
 \rho_e
 \mathbf w_\beta
 \dd V.
\]
The vector-potential block is
\[
 L_{\alpha\beta}
 =
 \mu
 \iint
 \mathbf w_\alpha(\mathbf r)^\ast
 G_k(\mathbf r,\mathbf r')
 \mathbf w_\beta(\mathbf r')
 \dd V'\dd V.
\]
For charge basis functions $\nu_r,\nu_s$,
\[
 \Phi_{rs}
 =
 \frac{1}{\epsilon}
 \iint
 \nu_r^\ast(\mathbf r)
 G_k(\mathbf r,\mathbf r')
 \nu_s(\mathbf r')
 \dd V'\dd V.
\]
The divergence block is
\[
 D_{r\alpha}
 =
 \int_{\Omega_c}
 \nu_r^\ast
 \nabla\cdot\mathbf w_\alpha
 \dd V.
\]
These four objects, together with terminal injection/extraction matrices, are
the minimal homogeneous-medium conductor teacher.

\subsection{Dimensionless assembly}

Before solving, scale
\[
 c=S_c\widetilde c,\qquad
 q=S_q\widetilde q
\]
using conductive-energy and charge-energy norms. The scaled block system is
assembled so resistance, inductive, and potential blocks have comparable
natural magnitude over the operating band. All Krylov residuals are evaluated
in a physical Riesz norm and then reported in dimensionless form.

\subsection{Nested enrichment policy}

The REFERENCE solver uses a deterministic nested sequence:
\[
 (p_s,r_\perp,q_{\rm quad})
 \rightarrow
 (p_s+\Delta p_s,r_\perp,q_{\rm quad})
 \rightarrow
 (p_s,r_\perp+\Delta r,q_{\rm quad})
 \rightarrow
 (p_s,r_\perp,q_{\rm quad}+\Delta q).
\]
Each direction is refined independently at least once so basis and quadrature
errors are identifiable.

For a port observable $Y$, define
\[
 \eta_s
 =
 \frac{\norm{Y_{s+}-Y}_{W_Y}}
      {\max(\norm{Y_{s+}}_{W_Y},Y_0)},
\]
and analogously $\eta_\perp,\eta_q$. A truth sample is accepted only if
\[
 \max(\eta_s,\eta_\perp,\eta_q,\eta_{\rm solve})
 \le\tau_{\rm ref}.
\]

\subsection{Reference assembly algorithm}

The first implementation follows this deterministic procedure:

\begin{enumerate}
\item Validate continuous conductor geometry and terminal topology.
\item Build Bishop frames and longitudinal quadrature partitions.
\item Build nested cross-section scalar, gradient, and curl-like modes.
\item Assemble the discrete divergence and exact terminal-current maps.
\item Remove charge gauge null spaces.
\item Assemble/cache intrinsic resistance and singular self-interaction blocks.
\item Assemble pair interactions using near/far classification.
\item Scale the mixed system and build a block preconditioner.
\item Solve all independent port excitations.
\item Evaluate true physical residual, reciprocity, passivity, and power closure.
\item Independently enrich longitudinal basis, cross-section basis, and
      quadrature.
\item Accept only when all declared convergence estimators pass.
\item Store the continuous/modal solution and certificate, never a world-grid
      field vector.
\end{enumerate}

\subsection{Junctions and branched conductors}

A junction joins multiple swept conductor branches at one electrical node.
Terminal-flux constraints are replaced by Kirchhoff compatibility:
\[
 \sum_{b\in\mathcal J}\mathcal I_b c_b=0
\]
for an internal junction. Potential/charge compatibility is imposed through a
shared junction charge/potential space or Lagrange multiplier. The first MVP
may exclude internal branched conductors, but the scene/port theory is designed
so adding them does not require a new global representation.

\subsection{Implementation checkpoint}

Before any neural network is trained, the homogeneous conductor teacher must
pass isolated-wire, single-loop, two-coil, rigid-motion, low-frequency, and
skin-depth convergence tests. A neural surrogate is not allowed to compensate
for a teacher that fails these tests.
